\documentclass[11pt,letterpaper]{article}
\usepackage[T1]{fontenc}
\usepackage[utf8]{inputenc}
\usepackage[margin=1in,headheight=14pt,headsep=18pt,footskip=28pt]{geometry}
\usepackage{newpxtext}
\usepackage{amsmath,amsthm}
\usepackage{newpxmath}
\usepackage{mathtools,microtype}
\usepackage{xcolor,booktabs,tabularx,array,enumitem}
\usepackage{fancyhdr,titlesec,needspace}
\usepackage[most]{tcolorbox}
\usepackage{xurl}
\definecolor{Navy}{HTML}{243D53}
\definecolor{Steel}{HTML}{3E617D}
\definecolor{Muted}{HTML}{636B72}
\definecolor{Pale}{HTML}{F2F5F7}
\usepackage[colorlinks=true,linkcolor=Steel,citecolor=Steel,urlcolor=Steel,
 bookmarksnumbered=true,pdfencoding=auto,psdextra]{hyperref}
\hypersetup{pdftitle={Every Countable Turing Ideal as Coarse Common Information},
 pdfauthor={Research manuscript prepared for Vladimir Reshetnikov},
 pdfsubject={Ideal realization, representative spectra, and fixed-core embeddings}}
\urlstyle{same}
\setlength{\parskip}{4pt plus 1pt minus .5pt}
\setlength{\parindent}{1.25em}
\setlength{\emergencystretch}{2em}
\linespread{1.025}
\setcounter{tocdepth}{1}
\titleformat{\section}{\Large\bfseries\color{Navy}}{\thesection}{.6em}{}
\titleformat{\subsection}{\large\bfseries\color{Navy}}{\thesubsection}{.6em}{}
\titlespacing*{\section}{0pt}{17pt}{8pt}
\titlespacing*{\subsection}{0pt}{12pt}{6pt}
\setlist{itemsep=3pt,topsep=5pt,parsep=0pt}
\pagestyle{fancy}
\fancyhf{}
\fancyhead[L]{\small\sffamily\color{Muted}COARSE COMMON INFORMATION}
\fancyhead[R]{\small\sffamily\color{Muted}IDEALS AND SPECTRA}
\fancyfoot[L]{\small\color{Muted}Research manuscript\quad\textbullet\quad 28 September 2026}
\fancyfoot[R]{\small\color{Muted}\thepage}
\renewcommand{\headrulewidth}{.3pt}
\fancypagestyle{plain}{\fancyhf{}\fancyfoot[R]{\small\thepage}\renewcommand{\headrulewidth}{0pt}}
\newtcolorbox{statusbox}{colback=Pale,colframe=Steel,boxrule=.45pt,
 arc=1pt,left=9pt,right=9pt,top=7pt,bottom=7pt,breakable}
\theoremstyle{plain}
\newtheorem{theorem}{Theorem}[section]
\newtheorem{lemma}[theorem]{Lemma}
\newtheorem{proposition}[theorem]{Proposition}
\newtheorem{corollary}[theorem]{Corollary}
\theoremstyle{definition}
\newtheorem{definition}[theorem]{Definition}
\newtheorem{example}[theorem]{Example}
\newtheorem{question}[theorem]{Question}
\theoremstyle{remark}
\newtheorem{remark}[theorem]{Remark}
\numberwithin{equation}{section}
\newcommand{\N}{\mathbb N}
\newcommand{\Cantor}{2^{\N}}
\newcommand{\leT}{\leq_{\mathrm T}}
\newcommand{\geT}{\geq_{\mathrm T}}
\newcommand{\nleT}{\nleq_{\mathrm T}}
\newcommand{\eqT}{\equiv_{\mathrm T}}
\newcommand{\ltT}{<_{\mathrm T}}
\newcommand{\leuc}{\leq_{\mathrm{uc}}}
\newcommand{\lenc}{\leq_{\mathrm{nc}}}
\newcommand{\equc}{\equiv_{\mathrm{uc}}}
\newcommand{\eqnc}{\equiv_{\mathrm{nc}}}
\newcommand{\Core}{\operatorname{Core}}
\newcommand{\CD}{\operatorname{CD}}
\newcommand{\Low}{\operatorname{Low}}
\newcommand{\Spec}{\operatorname{Spec}}
\newcommand{\Row}{\operatorname{Row}}
\newcommand{\Comp}{\mathrm{Comp}}
\newcommand{\degT}{\deg_{\mathrm T}}
\newcommand{\Rcode}{\mathcal R}
\newcommand{\Ecode}{\mathcal E}
\newcommand{\Jcode}{\mathcal J}
\newcommand{\sd}{\mathbin{\triangle}}
\newcommand{\rest}{\mathbin{\upharpoonright}}
\newcommand{\cat}{\mathbin{{}^\frown}}
\newcommand{\Dens}{\overline\rho}
\newcommand{\zero}{\mathbf 0}
\newcommand{\I}{\mathcal I}
\newcommand{\B}{\mathcal B}
\newcommand{\fiber}{\mathfrak F}
\newcommand{\upcone}[1]{\{\mathbf b:#1\leq\mathbf b\}}
\newcommand{\joinvec}{\bigoplus_{i\in\N}}
\newcommand{\patch}{\mathbin{\triangleright}}
\newcommand{\shortrepo}{\texttt{1a1396d4d3a2}}

\begin{document}
\thispagestyle{plain}
\begin{center}
{\sffamily\small\bfseries\color{Steel}COMPUTABILITY THEORY\quad /\quad RESEARCH CONTINUATION}\\[13pt]
{\fontsize{24}{29}\selectfont\bfseries\color{Navy}Every Countable Turing Ideal\\as Coarse Common Information}\\[10pt]
{\large Exact spectra and Turing-cone embeddings with a fixed core}\\[12pt]
{\color{Muted}Prepared for Vladimir Reshetnikov\\28 September 2026}
\end{center}
\vspace{6pt}
\begin{abstract}
For a set $X\subseteq\N$, let $\Core(X)$ be the sets computable from every
coarse description of $X$. We prove that the possible cores are exactly the
countable Turing ideals. More precisely, a positive-density column code for a
sequence $(A_i)$ has core equal to the ideal generated by that sequence, and its
entire representative spectrum is characterized by a single computable array
whose rows eventually stabilize as whole reals. The quantifier order in this
characterization is essential; merely computing each $A_i$ separately is not
enough.

We then separate common information from the difficulty of producing a
representative. If $Y$ can be approximated arbitrarily closely in upper-density
distance by computable sets, adjoining $Y$ preserves the core of every $X$.
For the classical dyadic code $\Rcode(H)$ this gives the exact operation
\[
 \Core(X\oplus\Rcode(H))=\Core(X),\qquad
 \Spec(X\oplus\Rcode(H))=\Spec(X)\cap\{\mathbf b:H\leT B'\}.
\]
Here $B$ is any representative of $\mathbf b$. Relative jump inversion turns
this identity into a join-preserving order embedding of the Turing cone above
$\degT(X')$ into both the uniform and the nonuniform coarse degrees with core
$\Core(X)$, above the coarse degree of $X$. Every fixed-core collection therefore
contains continuum many distinct representative spectra and has no countable
cofinal subset. The proofs are explicit, including the analytic locality lemma,
the coding estimates, and the jump-inversion argument.
\end{abstract}
\begin{statusbox}
\textbf{Scope and status.} This manuscript answers the core-realization question
left open in the inspected ProveIt research synthesis. Its coding ingredients
have substantial published antecedents, which are credited below. The paper
supplies proofs of the stated conclusions, but makes no claim of established
worldwide priority, independent refereeing, or Lean verification. In particular,
it does not resolve the general effective-dense least-representative question or
classify all possible coarse spectra.
\end{statusbox}
\newpage
\tableofcontents
\newpage

\section{The target and the contribution}\label{sec:intro}
\subsection{The repository question}
The ProveIt synthesis \cite{Repo} studies the failure of least Turing-degree
representatives in coarse equivalence classes. Its final subsection, ``What is
not settled'', specifically asks which countable ideals can occur as cores and
which upward-closed sets of degrees can occur as representative spectra. We
answer the first question completely and give a family of exact constructions
for the second. The repository was inspected at commit
\shortrepo; the full commit identifier and source path appear in
Appendix~\ref{app:sources}.

The distinction between these two questions is fundamental. The core records
what \emph{every} description can compute. The spectrum records which oracles
can produce \emph{some} description. A lower bound common to all descriptions
need not itself produce a description. A principal core is therefore not a
certificate of a least representative. The repository already demonstrates
this for every principal ideal \cite{Repo}; we do not count that observation as
a new result.

Gerdes's Question~7 \cite{Gerdes} places least representatives and spectral
classification in a broader comparison of robust reducibilities. The present
paper concerns binary coarse descriptions, where erroneous values are allowed
on a density-zero set. It does not silently transfer its constructions to
error-free partial descriptions.

\subsection{Four statements proved here}
The central results are the following, with all notation defined in the next
section.

\paragraph{Realization.}
For every countable Turing ideal $\I$ there is a binary set $X$ with
$\Core(X)=\I$. The construction works from any sequence of generators and gives
$X\eqT\joinvec A_i$ for that presentation
(Theorem~\ref{thm:realize}). Since every core is a countable ideal, this is an
exact classification of the possible cores.

\paragraph{Presentation-sensitive spectra.}
For our code $X_{\vec A}$, an oracle $B$ computes a coarse description exactly
when it computes a total binary array $f(i,s,n)$ satisfying
\[
 (\forall i)(\exists s_i)(\forall s\ge s_i)(\forall n)\quad
 f(i,s,n)=A_i(n).
\]
The stabilization is uniform over the entire real in each fixed row, but the
threshold need not be uniform over the rows (Theorem~\ref{thm:rows}). A concrete
example explains why the ideal alone cannot determine this spectrum.

\paragraph{Preserving the core while changing the spectrum.}
For arbitrary $X,H$, the operation $T_X(H)=X\oplus\Rcode(H)$ preserves
$\Core(X)$ and intersects its spectrum with the jump cone determined by $H$
(Theorem~\ref{thm:filter}). This is a special case of a general invariance theorem
for joins with sets of coarse computability bound $1$
(Theorem~\ref{thm:neutral}).

\paragraph{Universality within a fixed core.}
For $H,K\geT X'$,
\[
 H\leT K
 \quad\Longleftrightarrow\quad
 T_X(H)\leuc T_X(K)
 \quad\Longleftrightarrow\quad
 T_X(H)\lenc T_X(K).
\]
The induced maps preserve binary joins. Thus an entire Turing upper cone lives
above every coarse degree without changing its core
(Theorem~\ref{thm:embedding}). This is stronger than producing a pair of
inequivalent examples with the same common information.

\subsection{What comes from earlier work}
The cone-avoiding compactness theorem of Hirschfeldt, Jockusch, Kuyper, and
Schupp \cite[Theorem~3.7]{HJKS} is the source of the robust-radius principle
used here. The prefix-density proof in Section~\ref{sec:radius} follows the
locality approach in the repository synthesis \cite{Repo}. The classical dyadic
coarse-computability criterion is due to Jockusch and Schupp
\cite[Theorem~2.19]{JS}; we use its relativization and give the proof.

The positive-density column construction is not new. Hirschfeldt, Jockusch,
and Schupp \cite[Lemma~3.10 and Theorem~3.11]{HJS} use columns of robust codes of
the members of a countable ideal, together with finite-column truncations, in
studying closed subspaces of the density metric. Our exact core identity follows
by combining that mechanism with the radius principle. The repeated robust code
used for the sharper spectrum statement is a variant of the uniform coarsenings
in \cite[Section~2]{HJKS}; compare also the robust coding framework in
\cite{DI}.

Accordingly, the ideal-realization theorem should be understood as a precise
resolution of the repository's stated target from established ingredients,
not as evidence that a difficult published open problem had escaped all previous
methods. The whole-row characterization and fixed-core structural theorems are
proved as research extensions here. Their global novelty would require a more
exhaustive literature review than the targeted source check performed for this
manuscript.

\section{Descriptions, spectra, and common lower bounds}\label{sec:basic}
\subsection{Conventions}
We work in classical mathematics (ZFC) with standard oracle Turing machines.
We identify a set of natural numbers with its characteristic function. Write
$A\leT B$ for Turing reducibility, $A'$ for the Turing jump, and
$\degT(A)$ for the Turing degree of $A$. Binary join means
\[
 (A\oplus B)(2n)=A(n),\qquad (A\oplus B)(2n+1)=B(n).
\]
An infinite join uses a fixed computable pairing function. Different standard
pairings give uniformly Turing-equivalent encodings.

For $n\ge1$ and $E\subseteq\N$, put
\[
 \rho_n(E)=\frac{|E\cap[0,n)|}{n},\qquad
 \Dens(E)=\limsup_{n\to\infty}\rho_n(E),\qquad
 \delta(X,Y)=\Dens(X\sd Y).
\]
The pseudometric $\delta$ measures upper-density disagreement. We write
$X\approx Y$ when $\delta(X,Y)=0$, and
\[
 \CD(X)=\{D\in\Cantor:D\approx X\}.
\]
Every description in this paper is total and binary.

\begin{definition}
Write $X\lenc Y$ if every $D\in\CD(Y)$ computes some member of $\CD(X)$.
Write $X\leuc Y$ if a single Turing functional maps every member of $\CD(Y)$
to a member of $\CD(X)$. The corresponding equivalence relations are denoted
$\eqnc$ and $\equc$.

Set
\[
 \Spec(X)=\{\degT(B):(\exists D\leT B)\ D\in\CD(X)\},
 \qquad
 \Core(X)=\bigcap_{D\in\CD(X)}\Low(D),
\]
where $\Low(D)=\{A:A\leT D\}$.
\end{definition}

We use ideals of \emph{sets}: a Turing ideal is a nonempty collection closed
under Turing reducibility and binary join. It is automatically saturated under
Turing equivalence. Equivalently, one can work with its collection of degrees.
Countability has the same meaning in either formulation, because a Turing degree
contains only countably many sets.

\begin{lemma}[Elementary properties]\label{lem:corebasic}
For every $X$, $\Core(X)$ is a countable Turing ideal. It is the common lower
cone of $\Spec(X)$:
\[
 A\in\Core(X)\quad\Longleftrightarrow\quad
 (\forall\mathbf b\in\Spec(X))\ \degT(A)\le\mathbf b.
\]
Both $\Core$ and $\Spec$ are invariants of nonuniform coarse equivalence,
and hence of uniform coarse equivalence.
\end{lemma}
\begin{proof}
Downward closure is immediate. If two sets are computable from every
$D\in\CD(X)$, their binary join is too. Every computable set belongs to the
intersection. Since $X\in\CD(X)$, the core is contained in the countable
collection $\Low(X)$.

If $B$ computes a description of $X$, every core member is below $B$. Conversely,
each description itself supplies a degree in the spectrum. This proves the
common-lower-cone identity. Equivalence preserves which oracles can compute a
description, by composition of Turing computations, and then preserves the core.
\end{proof}

\begin{proposition}[The three representative spectra]\label{prop:spectra}
$\Spec(X)$ is upward closed and equals each of
\[
 \{\degT(D):D\approx X\},\quad
 \{\degT(Y):Y\equc X\},\quad
 \{\degT(Y):Y\eqnc X\}.
\]
Furthermore,
\begin{equation}\label{eq:order-spectrum}
 X\lenc Y\quad\Longleftrightarrow\quad\Spec(Y)\subseteq\Spec(X).
\end{equation}
\end{proposition}
\begin{proof}
Suppose $D\approx X$ and $D\leT B$. On the computable density-zero set
$P=\{2^n:n\in\N\}$ replace $D(2^n)$ by $B(n)$, and leave all other positions
unchanged. The resulting $E$ satisfies $E\approx X$ and $E\eqT B$: $B$ computes
$E$ using its computation of $D$, while the values on $P$ recover $B$.
This proves the first equality and upward closure. If $E\approx X$, the identity
functional gives reductions in both coarse directions, since
$\CD(E)=\CD(X)$.

If $Y\eqnc X$, applying $X\lenc Y$ to $Y$ itself shows that $Y$ computes a
description of $X$. This proves the remaining equalities. Finally, the forward
implication of \eqref{eq:order-spectrum} follows by composition. For the reverse,
a description $D$ of $Y$ has degree in $\Spec(Y)$, so the assumed inclusion
says exactly that $D$ computes a description of $X$.
\end{proof}

This proposition, including sparse degree coding, is background already present
in \cite{Repo}. The reversal of inclusion in \eqref{eq:order-spectrum} is worth
noting: a more difficult coarse problem has fewer solving degrees.

\begin{proposition}[Binary joins]\label{prop:joins}
For all $X,Y$,
\begin{equation}\label{eq:join-spectrum}
 \Spec(X\oplus Y)=\Spec(X)\cap\Spec(Y).
\end{equation}
Binary join is the least upper bound for both coarse reducibilities.
Moreover,
\begin{equation}\label{eq:join-distance}
 \delta(X\oplus Y,U\oplus V)
 \le\tfrac12\delta(X,U)+\tfrac12\delta(Y,V).
\end{equation}
The upper-density error of either half is at most twice the error of the join.
\end{proposition}
\begin{proof}
At a prefix of even length $2n$, the error count is the sum of the two length-$n$
error counts. Odd endpoints change the fractions by a quantity tending to zero.
This proves the density statements. Splitting a description gives descriptions
of both halves; conversely, joining two descriptions gives a description of the
join. These operations are uniform. The spectrum equality follows, allowing the
two reductions from a fixed solving oracle to have different indices.
\end{proof}

\begin{corollary}[Infimum versus least element]\label{cor:infimum}
$\Spec(X)$ has an infimum in the full partial order of Turing degrees if and
only if $\Core(X)$ is principal. If $\Core(X)=\Low(A)$, that infimum is
$\degT(A)$. It is a least element precisely when
$\degT(A)\in\Spec(X)$.
\end{corollary}
\begin{proof}
By Lemma~\ref{lem:corebasic}, the degrees of the core are exactly all the lower
bounds of the spectrum. They have a greatest member exactly when the ideal is
principal. A greatest lower bound is a least member precisely when it belongs
to the spectrum itself.
\end{proof}

\section{A locality theorem for common information}\label{sec:radius}
The following theorem is an equivalent form of the cone-avoiding compactness
principle in \cite[Theorem~3.7]{HJKS}. We include a proof to expose exactly where
nonuniform finite information enters. The prefix-metric approach is developed
in the repository's category section \cite{Repo}.

\begin{theorem}[Robust radius]\label{thm:radius}
For all $A,X\subseteq\N$,
\begin{equation}\label{eq:radius}
 A\in\Core(X)
 \quad\Longleftrightarrow\quad
 (\exists\eta>0)(\forall Y\in\Cantor)
 \bigl[\delta(X,Y)<\eta\ \Longrightarrow\ A\leT Y\bigr].
\end{equation}
\end{theorem}

\subsection{The complete metric used in the proof}
The pseudometric $\delta$ vanishes throughout $\CD(X)$, so it cannot be used for
Baire category on that space. Instead put
\[
 d(U,V)=\sup_{n\ge1}\rho_n(U\sd V).
\]
This is a genuine metric. It controls short prefixes as well as asymptotic
behavior, and $\delta(U,V)\le d(U,V)$.

\begin{lemma}\label{lem:metric}
The metric space $(\CD(X),d)$ is nonempty and complete. Every finite-prefix
cylinder is open in it. If $\delta(Y,Z)<a$, there is a finite modification
$\widetilde Y$ of $Y$ with $d(\widetilde Y,Z)<a$.
\end{lemma}
\begin{proof}
A disagreement at position $k$ forces $d(U,V)\ge1/(k+1)$, which proves both
separation and openness of cylinders. The triangle inequality follows from the
corresponding inequality for error counts at each prefix.

Let $(D_j)$ be $d$-Cauchy in $\CD(X)$. Each coordinate eventually stabilizes;
call the coordinatewise limit $D$. For a given $\varepsilon>0$, choose $j_0$ so
that $d(D_j,D_k)<\varepsilon/2$ whenever $j,k\ge j_0$. Fix $j\ge j_0$ and a
prefix length $n$. Taking $k$ sufficiently large that $D_k\rest n=D\rest n$
gives $\rho_n(D_j\sd D)<\varepsilon/2$. Consequently
$d(D_j,D)\le\varepsilon/2$, and $D_j\to D$ in $d$.
Since
\[
 \delta(D,X)\le d(D,D_j)+\delta(D_j,X)=d(D,D_j),
\]
we have $D\in\CD(X)$. Nonemptiness follows from $X\in\CD(X)$.

Finally choose $a'$ strictly between $\delta(Y,Z)$ and $a$. For all sufficiently
large $n$, $\rho_n(Y\sd Z)<a'$. Replace the corresponding finite initial
segment of $Y$ by that of $Z$. The new error count is zero at short prefixes
and no larger than the old error count at long prefixes. The resulting distance
is at most $a'<a$.
\end{proof}

The finite modification in the last paragraph is Turing equivalent to $Y$.
Its finite values and its length can be incorporated as constants into a
program. No effective procedure for finding those constants is asserted.

For a string $\tau$ and a set $Z$, let $\tau\patch Z$ agree with $\tau$ below
$|\tau|$ and with $Z$ afterwards. Then
\begin{equation}\label{eq:patch}
 d(\tau\patch Z,Z)=
 \max_{1\le m\le|\tau|}
 \frac{|\{j<m:\tau(j)\ne Z(j)\}|}{m},
\end{equation}
with the empty maximum interpreted as $0$. Beyond the end of the string the
numerator is fixed and the denominator only increases.

\subsection{Decoding from a neighborhood}
Fix a standard enumeration $(\Phi_e)$ of oracle programs, and define
\[
 S_{e,A}=\{D\in\CD(X):\Phi_e^D\text{ is total and equals }A\}.
\]
For a finite oracle string, a halting computation is accepted only when every
oracle query lies within that string.

\begin{lemma}[Local decoder]\label{lem:decoder}
Suppose $S_{e,A}$ is dense in a nonempty $d$-ball $B_d(Z,r)$ inside $\CD(X)$,
where $r>0$ is rational. Every $Y\in\Cantor$ satisfying $d(Y,Z)<r/4$ computes
$A$. A decoding program can be chosen using only $e$ and $r$, not the center $Z$.
\end{lemma}
\begin{proof}
On input $n$, search for a finite string $\tau$ such that $\Phi_e^\tau(n)$
halts and
\begin{equation}\label{eq:search}
 \max_{1\le m\le|\tau|}
 \frac{|\{j<m:\tau(j)\ne Y(j)\}|}{m}<r/2.
\end{equation}
Dovetail the halting tests. The displayed condition is a finite rational test
computable from $Y$. Output the value of the first successful computation.

The search terminates. Density supplies
$D_*\in S_{e,A}\cap B_d(Z,r/4)$, so $d(D_*,Y)<r/2$. A sufficiently long prefix
of $D_*$ witnesses the search condition and the halting computation at $n$.

Every accepted value is correct. For an accepted $\tau$, put
$H=\tau\patch Z$. It belongs to $\CD(X)$ because it is a finite modification
of $Z$. At prefixes of length at most $|\tau|$, the error against $Z$ is less
than $r/2+r/4$. By \eqref{eq:patch}, $H\in B_d(Z,r)$. The cylinder determined
by $\tau$, intersected with this ball, is therefore a nonempty open set.
It contains a member of $S_{e,A}$. On that member the same finite computation
must output $A(n)$, so the accepted output is $A(n)$.
\end{proof}

\begin{proof}[Proof of Theorem~\ref{thm:radius}]
The right-to-left direction holds because every coarse description has distance
zero from $X$. Conversely, let $A\in\Core(X)$. Then
$\CD(X)=\bigcup_e S_{e,A}$. The Baire category theorem and
Lemma~\ref{lem:metric} imply that some $S_{e,A}$ is not nowhere dense.
Thus its closure contains a nonempty ball $B_d(Z,r)$ with rational $r>0$,
and it is dense in that ball.

Let $\eta=r/4$, and suppose $\delta(X,Y)<\eta$. As $Z\approx X$,
$\delta(Y,Z)<\eta$. By the finite-patch part of Lemma~\ref{lem:metric}, a finite
modification $\widetilde Y\eqT Y$ satisfies $d(\widetilde Y,Z)<r/4$.
Lemma~\ref{lem:decoder} yields $A\leT\widetilde Y\eqT Y$.
\end{proof}

\begin{corollary}[Finite-stage capture]\label{cor:capture}
If $\delta(X,X_k)\to0$ and $X_k\leT B_k$ for every $k$, then
\[
 \Core(X)\subseteq\bigcup_k\Low(B_k).
\]
Indeed every member of the core is below $B_k$ for all sufficiently large $k$.
\end{corollary}
\begin{proof}
Choose the radius in Theorem~\ref{thm:radius} for a fixed core member and use
convergence. No effective rate and no uniform list of the reductions
$X_k\leT B_k$ are needed.
\end{proof}

This is the locality principle behind ideal realization: information common to
all exact coarse descriptions is already unavoidable at some positive error
scale. It can therefore be captured by a sufficiently accurate finite-stage
approximation.

\section{Coding geometry and a uniform robust code}\label{sec:codes}
\subsection{A partition with computably small tails}
Define
\begin{equation}\label{eq:pairing}
 p(i,m)=2^i(2m+1)-1,
 \qquad C_i=\{p(i,m):m\in\N\}.
\end{equation}
Every $n\in\N$ lies in exactly one column, whose index is $\nu_2(n+1)$.
Here $\nu_2(t)$ is the exponent of $2$ dividing the positive integer $t$.
The density of $C_i$ is $2^{-i-1}$.

\begin{lemma}[Column estimates]\label{lem:columns}
Let $T_k=\bigcup_{i\ge k}C_i$. Then
\begin{equation}\label{eq:tail}
 T_k=\{n:2^k\mid n+1\},\qquad
 |T_k\cap[0,N)|=\left\lfloor\frac{N}{2^k}\right\rfloor.
\end{equation}
If $E$ has density zero, then
$E_i=\{m:p(i,m)\in E\}$ has density zero for every $i$.
Conversely, if every $E_i$ has density zero, then $E$ has density zero.
In particular, a set meeting each column in only finitely many positions has
density zero.
\end{lemma}
\begin{proof}
The first assertions follow by factoring $n+1$. For the forward density
implication, the first $M$ positions of $C_i$ lie below $2^{i+1}M$, so
\[
 \frac{|E_i\cap[0,M)|}{M}
 \le 2^{i+1}\rho_{2^{i+1}M}(E)\longrightarrow0.
\]
For the reverse implication, fix $k$. On each of the finitely many columns
$i<k$, the count of errors below $N$ is $o(N)$, because the number of positions
of that column below $N$ is asymptotic to $2^{-i-1}N$. The remaining errors lie
in $T_k$. Hence
\[
 \limsup_N\rho_N(E)\le2^{-k}.
\]
Now let $k\to\infty$.
\end{proof}

The last argument is a finite-head/small-tail argument, not an appeal to closure
of density-zero sets under arbitrary countable unions. Such closure is false.
It is the uniform tail bound \eqref{eq:tail} that makes this particular
countable assembly safe.

\subsection{Three different codes}
We need to distinguish three operations. The dyadic repetition code is
\begin{equation}\label{eq:R}
 \Rcode(A)(p(i,m))=A(i).
\end{equation}
Its decoding from a coarse description generally costs a jump. In contrast,
write
\begin{equation}\label{eq:J}
 \Jcode(B)(0)=0,\qquad
 \Jcode(B)(t)=B(\lfloor\log_2 t\rfloor)\quad(t\ge1).
\end{equation}
Thus $B(k)$ is repeated on the entire interval $[2^k,2^{k+1})$.
Finally define
\begin{equation}\label{eq:E}
 \Ecode(A)=\Jcode(\Rcode(A)).
\end{equation}
All three exact codes are uniformly Turing equivalent to their inputs. The
composition in \eqref{eq:E} provides the stronger decoding property required
for our spectrum theorem.

\begin{lemma}[Block majority]\label{lem:majority}
For a binary oracle $D$, let $M_D(k)$ be the majority bit of $D$ on
$[2^k,2^{k+1})$, with ties assigned $0$. If $D\approx\Jcode(B)$, then
$M_D(k)=B(k)$ for all sufficiently large $k$.
\end{lemma}
\begin{proof}
A wrong majority on a block of length $2^k$ requires at least $2^{k-1}$ errors
on that block. Therefore
\[
 M_D(k)\ne B(k)\quad\Longrightarrow\quad
 \rho_{2^{k+1}}(D\sd\Jcode(B))\ge\tfrac14.
\]
The implication remains valid for $k=0$, where any wrong bit gives an even
larger error fraction. Infinitely many wrong block majorities would contradict
density-zero disagreement.
\end{proof}

\begin{proposition}[Whole-real recovery in the limit]\label{prop:robust}
There is a single oracle-computable procedure which, from any
$D\in\CD(\Ecode(A))$, produces a total binary function $F_D(s,n)$ with
\begin{equation}\label{eq:uniform-coarsening}
 (\exists s_0)(\forall s\ge s_0)(\forall n)\quad F_D(s,n)=A(n).
\end{equation}
In particular, every coarse description of $\Ecode(A)$ computes $A$.
\end{proposition}
\begin{proof}
Set
\[
 F_D(s,n)=M_D(p(n,s)).
\]
By Lemma~\ref{lem:majority}, there is a $K$ such that
$M_D(k)=\Rcode(A)(k)$ whenever $k\ge K$. Since $p(n,s)\ge s$ for every
$n,s$, all $s\ge K$ satisfy
\[
 F_D(s,n)=\Rcode(A)(p(n,s))=A(n)
 \quad\text{for every }n.
\]
This proves \eqref{eq:uniform-coarsening}. To obtain an individual Turing
reduction of $A$ to $D$, use one such stable $s$ as a fixed program constant.
\end{proof}

\begin{remark}[What is uniform, and what is not]\label{rem:uniform}
The procedure computing the entire two-variable array $F_D$ is uniform in $D$.
Selecting a stable row is not asserted to be uniform. Confusing these two claims
would invalidate later arguments. A single output real computed uniformly from
all descriptions is a stronger demand than the uniform array in
\eqref{eq:uniform-coarsening}.
\end{remark}

The repeated-code method is a concrete instance of uniform coarsening as used
in \cite[Definition~2.1 and Proposition~2.3]{HJKS}. The choice of exponentially
sized blocks rather than another standard rapidly growing partition affects
neither the conclusions nor the underlying computability notion.

\section{Realizing exactly the countable Turing ideals}\label{sec:ideals}
For a sequence $\vec A=(A_i)_{i\in\N}$, define
\begin{equation}\label{eq:idealcode}
 X_{\vec A}(p(i,m))=\Ecode(A_i)(m),
 \qquad
 \I_{\vec A}=\bigcup_{k\in\N}
 \Low\!\left(\bigoplus_{i<k}A_i\right).
\end{equation}
The empty finite join is computable. No computability assumption is placed on
the sequence of inputs; $\joinvec A_i$ is the oracle presenting that sequence.

\begin{theorem}[Ideal realization]\label{thm:realize}
For every sequence $\vec A$,
\begin{equation}\label{eq:realize}
 \Core(X_{\vec A})=\I_{\vec A},
 \qquad X_{\vec A}\eqT\joinvec A_i.
\end{equation}
Consequently the cores of binary sets are exactly the countable Turing ideals.
\end{theorem}
\begin{proof}
Let $D\approx X_{\vec A}$ and define the $i$th column oracle by
$D_i(m)=D(p(i,m))$. By Lemma~\ref{lem:columns},
$D_i\approx\Ecode(A_i)$. Proposition~\ref{prop:robust} gives $A_i\leT D_i\leT D$.
Thus every finite join of the $A_i$ is computable from $D$, and so is every
member of $\I_{\vec A}$. This proves
$\I_{\vec A}\subseteq\Core(X_{\vec A})$.

For the opposite inclusion let $X_{<k}$ agree with $X_{\vec A}$ on columns
$i<k$ and be zero on all later columns. Then
\[
 X_{<k}\leT\bigoplus_{i<k}A_i,
 \qquad
 \delta(X_{<k},X_{\vec A})\le2^{-k}
\]
by \eqref{eq:tail}. Corollary~\ref{cor:capture} now gives the desired reverse
inclusion. This is the step excluding information accidentally encoded across
infinitely many columns from becoming common to \emph{all} coarse descriptions.

The infinite join uniformly computes $X_{\vec A}$ by evaluating the codes.
In the other direction, $X_{\vec A}$ uniformly supplies each exact oracle
$\Ecode(A_i)$, and each $A_i(n)$ can be read at a computably specified position
of that exact code. For example,
\[
 A_i(n)=X_{\vec A}\bigl(p(i,2^{p(n,0)})\bigr).
\]
Hence the exact code computes the entire presentation, proving the Turing
 equivalence in \eqref{eq:realize}.

Finally let $\I$ be a countable Turing ideal and enumerate its members as
$(B_i)$. Take $A_i=\bigoplus_{j\le i}B_j$. These sets generate exactly $\I$,
so the construction realizes it. Conversely every core is a countable Turing
ideal by Lemma~\ref{lem:corebasic}.
\end{proof}

\begin{remark}[A shorter code suffices for the core alone]
Replacing $\Ecode(A_i)$ by $\Jcode(A_i)$ in \eqref{eq:idealcode} gives the same
core identity: block majority recovers $A_i$ modulo finitely many bits, and
those bits can be supplied as finite constants. This is closely aligned with
the coding construction in \cite[Theorem~3.11]{HJS}. We use $\Ecode$ because
its uniform approximation by whole reals also yields the exact spectrum in the
next section. The extra repetition is not needed to solve the bare
core-realization problem.
\end{remark}

\subsection{Two examples and an effectivity boundary}
\begin{example}[A principal ideal]
For the constant sequence $A_i=A$, the exact presentation is Turing equivalent
to $A$. Theorem~\ref{thm:realize} gives core $\Low(A)$; since the code itself
has degree $\degT(A)$, that degree is the least member of its spectrum.
\end{example}

\begin{example}[Exactly the arithmetical sets]\label{ex:arith}
Let $A_i=\emptyset^{(i)}$, the $i$th finite jump. Then
\[
 \Core(X_{\vec A})=
 \bigcup_i\Low(\emptyset^{(i)}),\qquad
 X_{\vec A}\eqT\emptyset^{(\omega)}
 :=\joinvec\emptyset^{(i)}.
\]
The displayed ideal is not principal. If it had a greatest degree, that degree
would be below some finite jump, whereas the next jump is a strictly larger
member of the ideal. Thus the representative spectrum of $X_{\vec A}$ has
\emph{no infimum at all}, by Corollary~\ref{cor:infimum}.
Every description still computes every individual finite jump. It need not
compute their infinite join: such a claim would make $\emptyset^{(\omega)}$
a core member, contradicting the formula above.
\end{example}

The construction is uniform from a \emph{given presentation}; it does not
provide a computable enumeration of an arbitrary countable ideal. Nor does
countability make the choices of indices for the separate reductions
$A_i\leT D$ uniformly computable. These distinctions are precisely what the
spectrum calculation measures.

\section{An exact spectrum theorem for the ideal code}\label{sec:rows}
\subsection{Eventually correct whole rows}
\begin{definition}\label{def:row}
For a fixed sequence $\vec A$, let $\Row(\vec A)$ be the set of Turing degrees
$\mathbf b$ for which an oracle $B\in\mathbf b$ computes a total binary
function $f(i,s,n)$ satisfying
\begin{equation}\label{eq:row}
 (\forall i)(\exists s_i)(\forall s\ge s_i)(\forall n)\quad
 f(i,s,n)=A_i(n).
\end{equation}
\end{definition}
The degree property does not depend on the chosen representative $B$.
The parameter $s_i$ may depend arbitrarily on $i$. However, once $s\ge s_i$,
the \emph{entire} function $n\mapsto f(i,s,n)$ is correct; this is not merely
pointwise convergence as $s$ grows.

\begin{theorem}[Whole-row spectrum]\label{thm:rows}
For the code in \eqref{eq:idealcode},
\begin{equation}\label{eq:rowspec}
 \Spec(X_{\vec A})=\Row(\vec A).
\end{equation}
This is simultaneously the spectrum of the density-zero class, the uniform
coarse class, and the nonuniform coarse class of $X_{\vec A}$.
\end{theorem}
\begin{proof}
Suppose $B$ computes $D\approx X_{\vec A}$. The column oracles
$D_i(m)=D(p(i,m))$ are uniformly computable from $B$ and satisfy
$D_i\approx\Ecode(A_i)$. Apply the single array-producing procedure of
Proposition~\ref{prop:robust} on each column:
\[
 f(i,s,n)=M_{D_i}(p(n,s)).
\]
For each fixed $i$, only finitely many block majorities are wrong. The proof
of that proposition supplies a threshold $s_i$ valid for every $n$, so
\eqref{eq:row} holds.

Conversely assume $B$ computes such an $f$. Construct a total binary set $D$
columnwise. Put $D(p(i,0))=0$. At $m\ge1$, put
\begin{equation}\label{eq:rowdecode}
 D(p(i,m))=f\!\left(i,m,
       \nu_2\bigl(\lfloor\log_2m\rfloor+1\bigr)\right).
\end{equation}
Indeed the target bit is
\[
 \Ecode(A_i)(m)
 =\Rcode(A_i)(\lfloor\log_2m\rfloor)
 =A_i\!\left(\nu_2\bigl(\lfloor\log_2m\rfloor+1\bigr)\right).
\]
For fixed $i$, \eqref{eq:rowdecode} is therefore correct at every
$m\ge\max(1,s_i)$, and is also correct at $m=0$ by definition of $\Ecode$.
The error set is finite on each column. Lemma~\ref{lem:columns} gives
$D\approx X_{\vec A}$, and $D\leT B$ by construction.
The last sentence follows from Proposition~\ref{prop:spectra}.
\end{proof}

The theorem is an exact effective condition rather than a characterization
solely in terms of the generated ideal. It removes the density quantifiers
from membership in this particular spectrum and replaces them by a uniform
array with the specified stabilization property.

\begin{corollary}[Necessary information bounds]\label{cor:rowbounds}
Let $S=\joinvec A_i$. Then
\begin{equation}\label{eq:sandwich}
 \{\mathbf b:S\leT B\}
 \ \subseteq\ \Spec(X_{\vec A})
 \ \subseteq\
 \{\mathbf b:(\forall i)\ A_i\leT B\ \text{and}\ S\leT B'\}.
\end{equation}
\end{corollary}
\begin{proof}
The left inclusion follows either from $X_{\vec A}\eqT S$ or by taking
$f(i,s,n)=A_i(n)$. For the right inclusion, a stable row computes $A_i$ from
$B$ using its threshold as a constant. Also $f(i,s,n)$ is a $B$-computable
pointwise convergent approximation to the characteristic function of $S$,
so the relativized limit lemma gives $S\leT B'$. A proof of that limit lemma
is included in Section~\ref{sec:jump}.
\end{proof}

\subsection{Computable rows can have a nontrivial collective spectrum}
\begin{example}[The presentation itself carries information]\label{ex:constantrows}
Given any $H\subseteq\N$, let $A_i$ be the constant real with value $H(i)$.
Every $A_i$ is computable, so $\I_{\vec A}=\Comp$. Nevertheless
\begin{equation}\label{eq:constantrows}
 X_{\vec A}\approx\Rcode(H),\qquad
 \Spec(X_{\vec A})=\{\mathbf b:H\leT B'\}.
\end{equation}
To see the first assertion, $\Ecode(A_i)$ agrees with the constant real $A_i$
at all positions except possibly $0$. Hence discrepancies from $\Rcode(H)$
can occur only at $p(i,0)=2^i-1$, a density-zero set. The spectral assertion
is Theorem~\ref{thm:jump-spectrum} below.

Taking $H=\emptyset''$ makes the spectrum omit the computable degree, even
though every individual input row is computable. Replacing this presentation
by the constant zero presentation gives the same ideal $\Comp$ but a spectrum
containing every Turing degree. Thus the condition
$(\forall i)\ A_i\leT B$ is not a substitute for \eqref{eq:row}.
\end{example}

This example is also a warning about the phrase ``uniformly computable
sequence''. A sequence all of whose entries are computable can encode an
arbitrarily complicated choice of those entries. Here that choice is exactly
$H$.

\section{The dyadic jump spectrum and uniform transport}\label{sec:jump}
This section gives the classical ingredients needed for the fixed-core
construction. They are not claimed as new: the basic dyadic criterion is the
relativization of \cite[Theorem~2.19]{JS}, and is also used in \cite{Repo,HJS}.

\begin{lemma}[Relativized limit lemma]\label{lem:limit}
$H\leT B'$ if and only if there is a total $B$-computable binary function
$h(i,s)$ such that $h(i,s)\to H(i)$ for every $i$.
\end{lemma}
\begin{proof}
For one direction, $B'$ can decide whether there exists a $t>s$ with
$h(i,t)\ne h(i,s)$, since this is a computably enumerable condition relative
to $B$. Search for a stage with no subsequent change and return its value.
The search terminates because the binary sequence converges.

For the other direction, approximate $B'$ by finite stages of its standard
$B$-computable enumeration. On input $(i,s)$, run the fixed reduction of $H(i)$
to $B'$ for $s$ steps, using that stage's finite approximation as the oracle.
Return its binary answer if it halts, and $0$ otherwise. The genuine halting
computation uses only finitely many oracle answers. Eventually all those
answers are correct and the time bound is long enough, so the simulation
returns $H(i)$ at all later stages.
\end{proof}

\begin{theorem}[Dyadic jump-cone criterion]\label{thm:jump-spectrum}
For every $H$,
\begin{equation}\label{eq:jumpspec}
 \Spec(\Rcode(H))=\{\mathbf b:H\leT B'\}.
\end{equation}
Moreover, for every $\varepsilon>0$ there is a computable $C$ such that
$\delta(\Rcode(H),C)<\varepsilon$.
\end{theorem}
\begin{proof}
Suppose $D\approx\Rcode(H)$ and $D\leT B$. On each column the errors have
density zero. Thus the majority of the first $s+1$ column values,
\[
 h(i,s)=\operatorname{Maj}\{D(p(i,0)),\ldots,D(p(i,s))\},
\]
is eventually $H(i)$. This is a total $B$-computable pointwise approximation,
so Lemma~\ref{lem:limit} gives $H\leT B'$.

Conversely, take the approximation $h$ supplied by that lemma and define
$D(p(i,m))=h(i,m)$. On each column $D$ differs from $\Rcode(H)$ at only
finitely many positions. The column lemma gives $D\approx\Rcode(H)$.

For the final claim, copy the finitely many bits $H(i)$ on columns $i<k$ and
put zero on columns $i\ge k$. The result $C_k$ is computable, because those
finitely many bits can be included as program constants, and
$\delta(C_k,\Rcode(H))\le2^{-k}$. Choose $k$ with $2^{-k}<\varepsilon$.
\end{proof}

The last assertion does not say that the sequence of indices for $C_k$ is
computable independently of $H$. It also does not imply the existence of a
single computable coarse description. For example, when $H=\emptyset''$,
\eqref{eq:jumpspec} rules out such a description.

\begin{proposition}[Uniform transport along Turing reductions]\label{prop:transport}
If $H\leT K$, then $\Rcode(H)\leuc\Rcode(K)$. In addition,
\begin{equation}\label{eq:Rjoin}
 \Rcode(H\oplus K)\equc\Rcode(H)\oplus\Rcode(K).
\end{equation}
\end{proposition}
\begin{proof}
Let $D\approx\Rcode(K)$. Column majorities give, uniformly in $D$, a total
array $k(i,s)$ with pointwise limit $K(i)$. Fix an index for the reduction
$H\leT K$. At stage $s$ simulate that program on input $i$ for $s$ steps with
oracle $j\mapsto k(j,s)$; return $0$ if no binary answer is obtained. The
resulting $h(i,s)$ is total and converges to $H(i)$, because the true computation
has a finite halting trace. Now output $h(i,m)$ at $p(i,m)$. The converse
construction in Theorem~\ref{thm:jump-spectrum} shows this is a coarse description
of $\Rcode(H)$. Every step is one fixed oracle procedure once the reduction
index is fixed.

For \eqref{eq:Rjoin}, a description of the right side splits into descriptions
of $\Rcode(H)$ and $\Rcode(K)$. Their column-majority approximations can be
interleaved into one approximation to $H\oplus K$, and then converted into a
description of $\Rcode(H\oplus K)$. Conversely, the even and odd entries of a
pointwise approximation to $H\oplus K$ approximate $H$ and $K$. Applying the
same conversion to both arrays and joining the outputs proves the other uniform
reduction.
\end{proof}

\section{Changing a spectrum without changing its core}\label{sec:neutral}
\subsection{A general invariance principle}
\begin{definition}
A set $Y$ is \emph{approximable by computable sets in density distance} if
\begin{equation}\label{eq:approximable}
 (\forall\varepsilon>0)(\exists C\text{ computable})\quad
 \delta(Y,C)<\varepsilon.
\end{equation}
Equivalently, its coarse computability bound
$\gamma(Y)=1-\inf_{C\text{ computable}}\delta(Y,C)$ equals $1$.
\end{definition}

The fact that \eqref{eq:approximable} implies $\Core(Y)=\Comp$ is the theorem
credited to Igusa in \cite[Theorem~4.3]{HJKS}. The following statement establishes
invariance under adjoining $Y$ to an \emph{arbitrary} base problem.

\begin{theorem}[Core neutrality under joins]\label{thm:neutral}
If $Y$ satisfies \eqref{eq:approximable}, then for every $X$,
\begin{equation}\label{eq:neutral}
 \Core(X\oplus Y)=\Core(X).
\end{equation}
\end{theorem}
\begin{proof}
A description of $X\oplus Y$ computes a description of $X$ by taking its even
half, so $\Core(X)\subseteq\Core(X\oplus Y)$.

For the reverse inclusion let $A\in\Core(X\oplus Y)$ and take its radius
$\eta>0$ from Theorem~\ref{thm:radius}. Choose a computable $C$ with
$\delta(Y,C)<\eta/2$. If $\delta(X,U)<\eta/2$, then
\[
 \delta(X\oplus Y,U\oplus C)
 \le\tfrac12\delta(X,U)+\tfrac12\delta(Y,C)<\eta.
\]
The radius property yields $A\leT U\oplus C\eqT U$.
Thus $\eta/2$ is a robust radius witnessing $A\in\Core(X)$.
\end{proof}

This proof uses more than $\Core(Y)=\Comp$. It uses arbitrarily accurate
\emph{computable approximations} to $Y$. The theorem does not assert that every
trivial-core set is neutral under joins.

\subsection{An exact jump filter}
\begin{definition}
For a base set $X$ and a parameter $H$, put
\[
 T_X(H)=X\oplus\Rcode(H),\qquad
 \mathcal Q_H=\{\degT(B):H\leT B'\}.
\]
\end{definition}

\begin{theorem}[Core-preserving spectral filter]\label{thm:filter}
For all sets $X,H$,
\begin{align}
 \Core(T_X(H))&=\Core(X),\label{eq:filtercore}\\
 \Spec(T_X(H))&=\Spec(X)\cap\mathcal Q_H.\label{eq:filterspec}
\end{align}
Also $X\leuc T_X(H)$ and $T_X(H)\eqT X\oplus H$.
For every sequence $\vec A$ this specializes to
\begin{equation}\label{eq:combined}
 \Core(T_{X_{\vec A}}(H))=\I_{\vec A},\qquad
 \Spec(T_{X_{\vec A}}(H))=\Row(\vec A)\cap\mathcal Q_H.
\end{equation}
\end{theorem}
\begin{proof}
Theorem~\ref{thm:jump-spectrum} shows that $\Rcode(H)$ satisfies
\eqref{eq:approximable}; hence \eqref{eq:filtercore} follows from
Theorem~\ref{thm:neutral}. Equations \eqref{eq:join-spectrum} and
\eqref{eq:jumpspec} give \eqref{eq:filterspec}. Projection to the even half
is the asserted uniform reduction. The exact Turing equivalence follows from
$\Rcode(H)\eqT H$. Finally use Theorems~\ref{thm:realize} and~\ref{thm:rows}.
\end{proof}

The parameter $H$ can be arbitrarily complicated. Its effect is visible in
which jumps occur above solving degrees, while the intersection of their lower
cones stays fixed. Thus the entire common-information ideal can remain
unchanged under substantial increases in the difficulty of producing a coarse
description.

\begin{corollary}[Strict extensions and unattained infima]\label{cor:strict}
If $H\nleT X'$, then
\[
 X\leuc T_X(H),\qquad T_X(H)\not\lenc X,
\]
and $\Spec(T_X(H))$ has no least element. If
$\Core(X)=\Low(A)$, its infimum is $\degT(A)$ and is not attained.
If $\Core(X)$ is nonprincipal, its spectrum has no infimum.
\end{corollary}
\begin{proof}
The degree of $X$ belongs to $\Spec(X)$ but, by
\eqref{eq:filterspec}, not to $\Spec(T_X(H))$. Formula
\eqref{eq:order-spectrum} rules out $T_X(H)\lenc X$.

Suppose $\mathbf b$ were the least member of $\Spec(T_X(H))$.
Every representative of $\mathbf b$ would lie in the core of $T_X(H)$,
which equals $\Core(X)\subseteq\Low(X)$. Thus $\mathbf b\le\degT(X)$.
Membership in the filtered spectrum would give
$H\leT B'\leT X'$, a contradiction. The two infimum assertions now follow
from Corollary~\ref{cor:infimum} and \eqref{eq:filtercore}.
\end{proof}

For example, $H=X''$ always satisfies the hypothesis. This gives a strict
extension inside the same core above every base set $X$, not merely above the
special ideal codes.

\begin{example}[A transparent principal-core spectrum]\label{ex:principal-filter}
Take $X=\Ecode(A)$. Proposition~\ref{prop:robust} and the exact equivalence
$X\eqT A$ give $\Spec(X)=\{\mathbf b:A\leT B\}$. Consequently
\[
 \Core(\Ecode(A)\oplus\Rcode(H))=\Low(A),\qquad
 \Spec(\Ecode(A)\oplus\Rcode(H))
 =\{\mathbf b:A\leT B\ \text{and}\ H\leT B'\}.
\]
When $H\nleT A'$, the infimum is $\degT(A)$ and is unattained. The existence
of prescribed unattained infima is already in \cite{Repo}; the displayed
spectral formula is the useful additional information supplied by this form of
the construction.
\end{example}

\section{Turing-cone universality inside a fixed core}\label{sec:embedding}
\subsection{The classical inversion ingredient}
We include the relative jump-inversion argument to make the separation part
of the embedding theorem independent of any unproved construction in the
repository. This is the classical Friedberg jump-inversion method, not a new
inversion theorem. Its use in density-metric arguments is also discussed in
\cite[Section~3]{HJS}.

\begin{lemma}[Relative jump inversion]\label{lem:inversion}
If $K\geT X'$, there is a set $B\geT X$ such that $B'\eqT K$.
\end{lemma}
\begin{proof}
We construct a binary real $Y$ and put $B=X\oplus Y$. Start with the empty
string $\sigma_0$. Given $\sigma_e$, ask whether there is a finite extension
$\tau\supseteq\sigma_e$ such that
\[
 \Phi_e^{X\oplus\tau}(e)\downarrow,
\]
where every query to the $Y$-part must be answered within $\tau$.
This is a computably enumerable question relative to $X$, hence decidable
using $X'$. If the answer is yes, choose the first witness found by a fixed
$X$-effective dovetailed search. If the answer is no, put $\tau=\sigma_e$.
In either case set
\[
 \sigma_{e+1}=\tau\cat K(e).
\]
Let $Y=\bigcup_e\sigma_e$. Each stage adds at least one bit, so $Y$ is total.
The construction is computable from $K$, because $K$ computes $X'$ and $X$.

The answer to the stage-$e$ question is exactly the characteristic value
$B'(e)$. A yes answer fixes a finite halting computation forever. A no answer
excludes halting on every infinite continuation of $\sigma_e$. Thus $K$
computes $B'$ by reconstructing the stages, and $B'\leT K$.

To decode $K$ from $Y\oplus X'$, reconstruct the same sequence of stages.
At stage $e$, $X'$ determines the yes/no answer and the deterministic choice
of $\tau$. Then the next coding bit is $Y(|\tau|)=K(e)$. Hence
$K\leT Y\oplus X'$. Since $Y\leT B\leT B'$ and $X'\leT B'$,
we obtain $K\leT B'$. The opposite reduction was already proved.
\end{proof}

The requirement $B\geT X$ matters. An arbitrary jump inverse of $K$ need not
solve the base coarse problem $X$; the relative construction guarantees that it
does.

\subsection{The embedding theorem}
For $r\in\{\mathrm{uc},\mathrm{nc}\}$ and a countable ideal $\I$, define the
\emph{fixed-core collection}
\[
 \fiber_r(\I)=\{[Y]_r:\Core(Y)=\I\},
\]
ordered by the corresponding coarse reducibility. It is well defined by
Lemma~\ref{lem:corebasic}. This is the fiber of the core invariant; no closure
under arbitrary joins of its members is assumed.

\begin{theorem}[Fixed-core upper-cone embedding]\label{thm:embedding}
Fix $X$ and let $\mathbf x=\degT(X)$. For $H,K\geT X'$,
\begin{equation}\label{eq:embed-equiv}
 H\leT K
 \quad\Longleftrightarrow\quad
 T_X(H)\leuc T_X(K)
 \quad\Longleftrightarrow\quad
 T_X(H)\lenc T_X(K).
\end{equation}
For either $r\in\{\mathrm{uc},\mathrm{nc}\}$, the map
\[
 \mathbf h\longmapsto[T_X(H)]_r
 \quad(\mathbf h\ge\mathbf x')
\]
is therefore a well-defined order embedding into $\fiber_r(\Core(X))$,
with every image element above $[X]_r$. It preserves binary joins:
\begin{equation}\label{eq:embedjoin}
 T_X(H\oplus K)\equc T_X(H)\oplus T_X(K).
\end{equation}
Distinct input degrees also have distinct representative spectra.
\end{theorem}
\begin{proof}
If $H\leT K$, split a description of $T_X(K)$ into its two halves. Keep its
$X$-description, and apply Proposition~\ref{prop:transport} to its description
of $\Rcode(K)$. Joining the outputs gives a description of $T_X(H)$ by one
fixed functional. This proves the first forward implication. Uniform
reducibility always implies nonuniform reducibility.

For the converse, suppose $T_X(H)\lenc T_X(K)$. By
Lemma~\ref{lem:inversion}, choose $B\geT X$ with $B'\eqT K$. The oracle $B$
computes the exact set $X$ and a coarse description of $\Rcode(K)$, the latter
by Theorem~\ref{thm:jump-spectrum}. Thus
\[
 \degT(B)\in\Spec(T_X(K)).
\]
The assumed nonuniform reduction and \eqref{eq:order-spectrum} put the same
degree in $\Spec(T_X(H))$. Formula \eqref{eq:filterspec} now gives
$H\leT B'\eqT K$, as required.

The map is well defined on Turing degrees by the already proved forward
implication in both directions, and is injective by the converse. Its core and
its position above $[X]_r$ are given by Theorem~\ref{thm:filter}.

For binary joins, rearrange
\[
 (X\oplus\Rcode(H))\oplus(X\oplus\Rcode(K))
\]
to two copies of an $X$-description and one description of each dyadic code.
Either copy of the $X$-description suffices; in the reverse direction one can
duplicate it. These finite splitting, rearranging, and duplicating operations
preserve density-zero error. Proposition~\ref{prop:transport} replaces the two
dyadic descriptions by one of $\Rcode(H\oplus K)$, or conversely.
This proves \eqref{eq:embedjoin} with uniform reductions.

Finally, equality of two image spectra implies nonuniform equivalence by
\eqref{eq:order-spectrum}, and \eqref{eq:embed-equiv} then forces equality of
the input Turing degrees.
\end{proof}

The lower endpoint $\mathbf x'$ of the domain is not claimed to map to
$[X]_r$. The theorem is an embedding into the region above $[X]_r$, not a
bottom-preserving isomorphism with that entire region. Likewise it identifies a
join-closed \emph{subcollection} of the fiber; it does not show that the full
fiber is closed under arbitrary binary coarse joins.

\subsection{Consequences for all countable ideals}
\begin{corollary}[Continuum many spectra with one prescribed core]\label{cor:continuum}
For every countable Turing ideal $\I$ and either coarse reducibility,
$\fiber_r(\I)$ contains a join-preserving copy of a Turing upper cone. Above
each of its members there are continuum many members with pairwise distinct
representative spectra.
\end{corollary}
\begin{proof}
Theorem~\ref{thm:realize} makes the fiber nonempty. Apply
Theorem~\ref{thm:embedding} to any representative $X$ of a chosen member.
Every Turing upper cone has continuum many degrees: there are continuum many
oracles of the form $X'\oplus Z$, and each Turing degree contains only
countably many oracles. Distinctness of the image spectra is part of that theorem.
\end{proof}

\begin{corollary}[No maximal member and no countable cofinal family]\label{cor:cofinal}
For every countable Turing ideal $\I$ and
$r\in\{\mathrm{uc},\mathrm{nc}\}$, the ordered collection
$\fiber_r(\I)$ has no maximal member and no countable cofinal subset.
\end{corollary}
\begin{proof}
For any $X$ with core $\I$, the set $T_X(X'')$ has the same core and is
strictly above $X$ even for nonuniform coarse reducibility, by
Corollary~\ref{cor:strict}. This proves the first assertion.

For the second, let $(X_n)$ be any nonempty countable sequence of representatives
with core $\I$, repeating entries if the proposed family is finite. Put
\[
 L=\bigoplus_n X_n',\qquad H=L',\qquad Z=T_{X_0}(H).
\]
Then $\Core(Z)=\I$. If $Z\lenc X_n$, the exact oracle $X_n$ would compute a
coarse description of $Z$. Formula \eqref{eq:filterspec} would imply
$H\leT X_n'\leT L$, contradicting $L'\nleT L$.
Thus $Z$ is not below any listed member even nonuniformly, so the family is
not cofinal for either reducibility. The empty family is not cofinal because
the fiber is nonempty.
\end{proof}

Cofinality here is inside the \emph{coarse-degree collection with a fixed
core}. It must not be confused with the coinitiality of Turing degrees inside
one representative spectrum, which is a different question studied in
\cite{Repo}.

\section{Interpretation, limitations, and formalization}\label{sec:interpret}
\subsection{What has been determined}
There is no additional realizability restriction on a countable Turing ideal
arising from being the information common to all coarse descriptions. Every
such ideal occurs, with a code whose exact degree is the degree of its chosen
presentation. In particular, cores need not be principal, and the lack of a
least representative can reflect either an unattained infimum or the complete
absence of an infimum.

At the same time, this classification is far from a classification of coarse
degrees. Theorem~\ref{thm:embedding} locates a full Turing upper cone inside
every one of the classes grouped together by the core invariant. Those embedded
problems can be ordered, joined, and separated by their spectra while preserving
every member of their common lower cone. The failure of completeness of the
core invariant is therefore structural, not just an isolated counterexample.

The construction also separates two kinds of uniformity. A coarse description
of $X_{\vec A}$ uniformly produces an array approximating all its rows. It
need not uniformly identify the stabilized row for each $i$. That missing
selection information is a genuine source of spectral complexity, already
visible when every individual $A_i$ is computable.

\subsection{Claims not made}
The paper does not show that every upward-closed collection of Turing degrees
with a prescribed common lower cone is a coarse spectrum. Formulas
\eqref{eq:rowspec} and \eqref{eq:combined} describe the particular constructions
used here. They do not imply that every possible spectrum has one of these forms.

A spectrum without a least member can still have minimal members. None of the
arguments proves the absence of all minimal elements in a representative
spectrum. Similarly, the upper-cone embedding uses no claim that Turing jumps
preserve arbitrary meets or that the Turing degrees form a lattice.

No computable rate of convergence is supplied for arbitrary sequences of
generators, and no arithmetical complexity bound follows from countability
alone. The explicit bound $X_{\vec A}\eqT\joinvec A_i$ is relative to a selected
presentation, not an intrinsic canonical degree assigned to the ideal.

Finally, density-zero modifications preserve coarse descriptions because errors
are permitted. Effective-dense descriptions permit omissions but not erroneous
answers. Sparse replacement and the coding constructions cannot be reused for
that setting without a separate proof. Gerdes's effective-dense question
\cite[Question~7]{Gerdes}, also marked C2 in \cite{Repo}, is not answered here.

\subsection{A formalization route}
The results have been proved as conventional mathematics in this manuscript;
no accompanying Lean theorem file has been compiled. A useful formalization
order would isolate three modules.

First, formalize finite density estimates and the partition $p(i,m)$:
the exact tail count, restriction of null errors to a positive-density column,
the finite-head/small-tail assembly lemma, and block-majority recovery. These
are the parts directly supported by the finite checks in this package, although
finite checks do not establish their universal statements.

Second, formalize the common-information invariant and its robust-radius theorem.
This requires the prefix-density metric, its completeness on $\CD(X)$,
continuity of finite oracle computations, and the local decoder. The dependence
on Baire category should be exposed rather than hidden in an admitted
``compactness'' statement. With that module available, the ideal-realization
and neutrality theorems are short deductions.

Third, add use-bounded oracle computations, the limit lemma, and relative jump
inversion. The proof of Theorem~\ref{thm:embedding} then separates into a uniform
positive reduction and a single inverse-jump witness for the converse. A clean
API should distinguish a Turing reduction of an individual real from an oracle
procedure returning a convergent array. That distinction is essential in
Theorem~\ref{thm:rows}.

\section{Further research questions}\label{sec:questions}
The questions below are proposed continuations of the proved results. They are
not all asserted to be independently verified open problems in the published
literature. Each identifies a boundary not settled by this manuscript.

\begin{question}[Spectra over a fixed ideal]\label{q:spectra}
Fix a countable ideal $\I$. Which upward-closed sets $S$ of Turing degrees occur
as $\Spec(X)$ with $\Core(X)=\I$? Beyond the common-lower-cone condition, what
effective or descriptive-set-theoretic restrictions are necessary and sufficient?
\end{question}
Theorem~\ref{thm:filter} gives exact closure operations on already realized
spectra. A first concrete target is to find a necessary condition not implied
merely by upward closure and the prescribed common lower cone.

\begin{question}[Simplifying the whole-row condition]\label{q:rows}
For which presentations $\vec A$ is the right inclusion in
\eqref{eq:sandwich} an equality? Can whole-row stabilization be characterized
by a useful index-selection or domination condition on a presentation?
\end{question}
The constant-row example rules out dropping the jump condition altogether.
It does not settle the converse with both necessary conditions present.
Canonical increasing sequences of finite jumps provide concrete test cases.

\begin{question}[Exact characterization of neutral joins]\label{q:neutral}
Is a set $Y$ neutral for all cores, in the sense that
$\Core(X\oplus Y)=\Core(X)$ for every $X$, exactly when $\gamma(Y)=1$?
If not, what is the correct characterization?
\end{question}
Only the sufficient direction is proved here. Taking $X$ computable shows
that universal neutrality implies a trivial core, but this necessary condition
is weaker than the hypothesis used in Theorem~\ref{thm:neutral}.

\begin{question}[When common information increases under joins]\label{q:synergy}
Given $\Core(X)=\Core(Y)=\I$, under what additional assumptions does
$\Core(X\oplus Y)=\I$ follow? Which larger ideals can arise from their join?
\end{question}
The embedded families in Theorem~\ref{thm:embedding} stay in one core under
binary joins. This does not settle the behavior of arbitrary pairs from the
same fixed-core collection.

\begin{question}[Uniform degree structure beyond spectra]\label{q:uniform}
Within a fixed core, how much can uniform coarse degrees vary while the
nonuniform coarse degree, and therefore the representative spectrum, is fixed?
Can one give an exact presentation-theoretic criterion for that variation?
\end{question}
Equation~\eqref{eq:order-spectrum} completely captures the nonuniform order,
but not the requirement for a single functional. The explicit uniform
reductions in Theorem~\ref{thm:embedding} do not collapse the two global degree
structures.

\begin{question}[Minimal solving degrees]\label{q:minimal}
For the spectra $\Row(\vec A)\cap\mathcal Q_H$, when do minimal elements
exist, and can one give a natural presentation whose spectrum has none?
\end{question}
The existence of an infimum is already decided by the core. Minimality is a
separate local property of the upward-closed spectrum and calls for a different
argument, likely involving controlled jump inversions or refinements of
presentations.

\begin{question}[Effective realizations and formal assurance]\label{q:effective}
For ideals presented at specified arithmetical or hyperarithmetical complexity,
what are the sharp degrees of realizing codes and of cone-avoiding descriptions?
Which parts of the radius and whole-row theorems admit economical Lean
formalizations over the repository's existing oracle-computability framework?
\end{question}
An effective index for a sequence is substantially more information than a
nonuniform collection of individually simple generators. Sharp results should
state explicitly which presentation data are supplied to the construction.

\begin{question}[The error-free boundary]\label{q:ed}
What replaces the core and the whole-row spectrum theorem for effective-dense
descriptions, and can an error-free coding operation preserve their common
information while imposing an independent jump condition?
\end{question}
A successful answer would have to address the prohibition on erroneous values,
not merely shrink their density. It could provide a genuinely new route to the
remaining effective-dense least-representative problem.

\appendix
\section{Proof audit and quantitative checks}\label{app:audit}
\subsection{Critical logical checkpoints}
The following checkpoints summarize the issues most likely to create a false
stronger statement.

\paragraph{Finite constants are not a uniform algorithm.}
The finite patch in Theorem~\ref{thm:radius}, the stable row in
Proposition~\ref{prop:robust}, and the finite truncations in
Theorem~\ref{thm:jump-spectrum} are used as witnesses to individual Turing
reductions. The paper never claims that their indices or stabilization stages
are computable uniformly from the displayed oracles.

\paragraph{Whole-row stabilization is stronger than pointwise convergence.}
The decoding formula \eqref{eq:rowdecode} uses the value at a coordinate that
itself changes with the stage $m$. It is justified by
$\exists s_i\,\forall s\ge s_i\,\forall n$, not by
$\forall n\,\exists s_{i,n}\,\forall s\ge s_{i,n}$.
The block-of-repetitions code was chosen specifically to supply the stronger
order of quantifiers.

\paragraph{Countably many error sets require the tail estimate.}
Nullity of every column's error set does not by itself follow from a countable
union theorem. The proof first handles finitely many columns, then bounds
\emph{all} remaining errors by $2^{-k}$, and only then lets $k$ grow.
The exact floor formula in \eqref{eq:tail} is the quantitative certificate.

\paragraph{Core realization needs both inclusions.}
Robust recovery of each generator proves only
$\I_{\vec A}\subseteq\Core(X_{\vec A})$. The upper inclusion is a separate
application of robust radius to the finite-column truncations. Omitting it
would leave open the possibility of unintended common information across the
columns.

\paragraph{The reverse embedding uses a relative jump inverse.}
The witness $B$ must satisfy both $B\geT X$ and $B'\eqT K$.
These properties ensure that $B$ actually solves the target base-plus-code
problem before its jump separates $H$ from $K$. Lemma~\ref{lem:inversion}
proves exactly this pair of requirements.

\paragraph{Different partial orders are kept separate.}
The notions of least member, minimal member, infimum, cofinal family, and
coinitial family are not interchanged. In particular,
Corollary~\ref{cor:cofinal} concerns coarse degrees in a fixed-core collection;
Corollary~\ref{cor:infimum} concerns Turing degrees within one spectrum.

\subsection{Finite verification artifacts}
The companion program \texttt{checks/finite\_checks.py} uses only Python's
standard library. It checks the column partition and exact tail counts,
majority-error thresholds, the index bound used for whole-real stabilization,
and finite simulations of the assembly and presentation examples. It also
checks the finite-prefix triangle and patch bounds used in the local decoder.
A machine-readable result is provided as \texttt{checks/results.json}.

These are finite combinatorial checks. They do \emph{not} verify the Baire
category theorem, Turing reducibility claims, convergence at unbounded stages,
jump inversion, global novelty, or any full theorem of this paper.
The mathematical assurance for those assertions is the written proof, pending
independent review or formalization.

\section{Source ledger and reproducibility}\label{app:sources}
The source-level comparison used the following repository snapshot:
\begin{quote}\small
\textbf{Repository:} \url{https://github.com/VladimirReshetnikov/ProveIt}\\
\textbf{Commit:} \texttt{1a1396d4d3a2ac6812692df520517d86ba3a4785}\\
\textbf{Primary research source:}
\path{Computability/TuringDegrees/Research/CoarseDegrees/research-synthesis/Turing_Degrees_Synthesis.tex}
\end{quote}
The relevant source locations are its section on spectra, its category section
on robust radius, and its final subsection on unsettled questions. In particular,
the classification question is explicitly present in that final subsection.
The repository is a research source, not a substitute for independent verification.
All ingredients used here have been proved in this manuscript, apart from
standard basic facts about oracle computation, the Turing jump, and Baire
category.

The targeted literature check covered the original coarse-reducibility paper
\cite{HJKS}, the dyadic criterion \cite{JS}, the countable-column antecedent
\cite{HJS}, the robust-coding framework \cite{DI}, and the motivating question
\cite{Gerdes}. The comparison does not certify absence of these conclusions
from all other papers, preprints, theses, or unpublished notes.

The archive includes the self-contained source \texttt{article.tex}, the compiled
\texttt{article.pdf}, a three-pass \texttt{build.sh}, a concise proof-status file,
a source ledger, and the finite checks. The bibliography is embedded in the TeX
source; no external \texttt{.bib} file or network access is needed to rebuild.
The build requires a standard \LaTeX{} installation with the packages listed
in the source. No Lean verification is included or implied.

\begingroup\small\raggedright
\begin{thebibliography}{9}
\bibitem{Repo}
Vladimir Reshetnikov's ProveIt repository.
\emph{Coarse Classes Without Least Turing Degrees: Deduplicated Results of the
Nine Research Reports on Target C1}. Research synthesis, September 2026.
Inspected at commit \shortrepo, 28 September 2026 (Pacific time).
Source path and full revision identifier in Appendix~\ref{app:sources}.
\url{https://github.com/VladimirReshetnikov/ProveIt/tree/1a1396d4d3a2ac6812692df520517d86ba3a4785/Computability/TuringDegrees/Research/CoarseDegrees}.

\bibitem{HJKS}
Denis R. Hirschfeldt, Carl G. Jockusch, Jr., Rutger Kuyper, and Paul E. Schupp.
Coarse reducibility and algorithmic randomness.
\emph{Journal of Symbolic Logic} \textbf{81} (2016), 1028--1046.
\url{https://doi.org/10.1017/jsl.2015.70}.
Preprint: \url{https://arxiv.org/abs/1505.01707}.

\bibitem{JS}
Carl G. Jockusch, Jr., and Paul E. Schupp.
Generic computability, Turing degrees, and asymptotic density.
\emph{Journal of the London Mathematical Society} \textbf{85}(2) (2012), 472--490.
\url{https://doi.org/10.1112/jlms/jdr051}.
Preprint: \url{https://arxiv.org/abs/1010.5212}.

\bibitem{HJS}
Denis R. Hirschfeldt, Carl G. Jockusch, Jr., and Paul E. Schupp.
\emph{Coarse computability, the density metric, Hausdorff distances between
Turing degrees, perfect trees, and reverse mathematics}.
Preprint, 2021. References to Lemma~3.10 and Theorem~3.11 use the inspected
arXiv text.
\url{https://arxiv.org/abs/2106.13118}.

\bibitem{DI}
Damir D. Dzhafarov and Gregory Igusa.
Notions of robust information coding.
\emph{Computability} \textbf{6}(2) (2017), 105--124.
\url{https://doi.org/10.3233/COM-160059}.
Preprint: \url{https://arxiv.org/abs/1406.2982}.

\bibitem{Gerdes}
Peter M. Gerdes.
\emph{Comparing Notions of Dense Computability on $\omega^\omega$ and
$2^\omega$}. arXiv:2508.06925v1, 2025. Question~7.
\url{https://arxiv.org/abs/2508.06925v1}.
\end{thebibliography}
\endgroup
\end{document}
